\documentclass{article} %210 mm × 297 mm
\usepackage[utf8]{inputenc}
\usepackage[a4paper,left=1.5cm, right=1.5cm, top=2cm, bottom=2cm]{geometry}
\usepackage{graphicx}
%\usepackage{blindtext}
\usepackage{multirow}
\usepackage{mhchem}
\usepackage{url}
\usepackage{subfigure}
\usepackage{amsfonts,latexsym} % para tener disponibilidad de diversos simbolos
\usepackage{enumerate}
%\usepackage{booktabs}
\usepackage{float}
%\usepackage{threeparttable}
\usepackage{array,colortbl}
%\usepackage{ifpdf}
%\usepackage{rotating}
\usepackage{cite}
\usepackage{stfloats}
\usepackage{url}
\usepackage{listings}
\usepackage{fancyhdr}

\usepackage{color}
\usepackage{hyperref}
%\usepackage{wrapfig}
\usepackage{array}
%\usepackage{multirow}
%\usepackage{adjustbox}
%\usepackage{nccmath}
\usepackage[dvipsnames]{xcolor}
\usepackage{todo}

\newcommand{\titulo}{Anleitungen und Formeln MfN 2}
\newcommand{\FCE}{\textbf{SoSe 2025}}
\newcommand{\R}[0]{\mathbb{R}}
\newcommand{\Ra}[0]{\Rightarrow}
\newcommand{\ram}[0]{$\Ra$ }
\renewcommand{\xi}[0]{\textcolor{CornflowerBlue}{\mathbf{x_1}}}
\newcommand{\xii}[0]{\textcolor{OliveGreen}{\mathbf{x_2}}}
\newcommand{\xn}[0]{\textcolor{BrickRed}{\mathbf{x_n}}}

\usepackage{fancyhdr}
\pagestyle{fancy}
\fancyhead{}
\fancyfoot{}
\fancyhead[LO]{\small\nombre}
\fancyhead[RE]{\small\titulo}
\fancyhead[LO]{\small\FCE}
\fancyfoot[RO,LE] {\thepage}
 
\setcounter{page}{1} %Número de la página, la editorial lo cambiará



\usepackage[onehalfspacing]{setspace}
\begin{document}
\tableofcontents
\section{Grundlagen}
\begin{equation}
    u = \ln(x) \quad \Rightarrow \quad x = e^u
\end{equation}

\begin{equation}
    (\cos{x})^2 + (\sin{x})^2 = 1
\end{equation}
Implikationen: 
\begin{align*}
    (\sin{x})^2 & = 1 - (\cos{x})^2 \\
    - (\sin{x})^2 & = (\cos{x})^2 -1 \\
    (\cos{x})^2 & = 1- (\sin{x})^2 \\
   - (\cos{x})^2 & = (\sin{x})^2 -1 \\
\end{align*}

\section{Definitheit}
\begin{table}[H]
    \centering
    \begin{tabular}{l|l}
        alle EW $>$ 0 & positiv definit \\
        alle EW $\geq$ 0 & positiv semidifinit \\
        alle EW $<$ 0 & negativ definit \\
        alle EW $\leq$ 0 & negativ semidifinit \\
        EW $<$ 0 und $>$ 0 & indefinit
    \end{tabular}
    \caption{Definitheit}
    \label{tab:Definitheit}
\end{table}

\section{Ableitungen}
\subsection{Sollte man kennen}
\begin{align}
    \frac{d}{dx} ln(x) &  = \frac{1}{x} \\
    \frac{d}{dx} \sin (x) & = \cos(x) \\
    \frac{d}{dx} \cos (x) & = -\sin(x) \\
\end{align}
\subsection{Partielle Ableitung} \label{subsec:partial}
Eine Funktion mit mehreren Variablen nach nur einer ableiten. 
\begin{equation}
    \frac{\partial f}{\partial x} f(x,y) =  \frac{\partial f}{\partial x}
\end{equation}
\subsection{Gradient} \label{subsec:gradient}
\begin{equation}
    grad(f)(x,y) = \Big(\frac{\partial}{\partial x} f(x,y),  \frac{\partial}{\partial y} f(x,y)\Big)
\end{equation}
Also
\begin{equation*}
    grad(f)(x,y) = \Big(\text{PA nach x},  \text{PA nach y}\Big)
\end{equation*}
\subsection{Hesse-Matrix} \label{subsec:hessematrix}
\begin{equation}
    H(f)(x,y) = \begin{pmatrix}
        \frac{\partial}{\partial x}   \frac{\partial}{\partial x} f(x,y) & \frac{\partial}{\partial x}   \frac{\partial}{\partial y} f(x,y) \\
        \frac{\partial}{\partial y}   \frac{\partial}{\partial x} f(x,y) & 
        \frac{\partial}{\partial y}   \frac{\partial}{\partial y} f(x,y)
    \end{pmatrix}
\end{equation}
Also 
\begin{equation*}
    H(f)(x,y) = \begin{pmatrix}
        \text{PA nach x, dann erneut nach x } & \text{PA nach x, dann nach y} \\
       \text{PA nach y, dann nach x}  & 
        \text{PA nach y, dann erneut nach y}
    \end{pmatrix}
\end{equation*}
Oder auch das Muster: 
\begin{equation*}
    H(f)(x,y) = \begin{pmatrix}
        \text{PA nach } \xi (\text{PA nach } \xi ) & \text{PA nach } \xi (\text{PA nach } \xii ) & \cdots & \text{PA nach } \xi (\text{PA nach } \xn ) \\
        \text{PA nach } \xii (\text{PA nach } \xi ) & \text{PA nach } \xii (\text{PA nach } \xii ) & \cdots & \text{PA nach } \xii (\text{PA nach } \xn ) \\
        \vdots & & & \vdots \\
      \text{PA nach } \xn (\text{PA nach } \xi ) & \text{PA nach } \xn (\text{PA nach } \xii ) & \cdots & \text{PA nach } \xn (\text{PA nach } \xn ) \\
    \end{pmatrix}
\end{equation*}

\section{Extrema bestimmen}
\paragraph{Für $f: \R \times \R \rightarrow \R$}
\begin{enumerate}
    \item Bestimme den \nameref{subsec:gradient} $grad(f)(x,y)$
    \item Finde alle Kandidaten: Finde alle $(x,y)$ mit $grad(x,y) = (0,0)$
    \item Bestimme die \nameref{subsec:hessematrix} $H(f)(x,y)$ für alle Kandidaten
    \begin{itemize}
        \item $H$ positiv definit $\Ra$ Minimum bei $(x,y)$
        \item $H$ negativ definit: $\Ra$ Maximum bei $(x,y)$
        \item $H$ indefinit: $\Ra$ Sattelpunkt bei $(x,y)$ \\ Definitheit siehe Tabelle \ref{tab:Definitheit} %\nameref{tab:Definitheit}
    \end{itemize}
\end{enumerate}
\subsection{Hurwitzkriterium}
\paragraph{Für 2 $\times$ 2 Matrizen}
\[
A = \begin{pmatrix}
a & b \\  c & d
\end{pmatrix}
\]
\begin{enumerate}
    \item $det(A) < 0$ \ram indefinit
    \item $det(A) > 0$: 
    \begin{enumerate}
        \item $a < 0$ \ram negativ definit
        \item $a > 0$ \ram positiv definit
    \end{enumerate}
\end{enumerate}

\paragraph{Für 3 $\times$ 3 Matrizen}
\[
A  = \begin{pmatrix}
a & b & c\\   d & e & f \\ g & h & i
\end{pmatrix}
\]
\begin{table}[H]
    \centering
    \begin{tabular}{l|l|l|ll}
       $a$  & $det \begin{pmatrix}a&b \\ d & e \end{pmatrix}$  & $det(A)$ & \ram    \\ \hline
       $> 0$  & $ > 0$ & $>0$ & positiv definit & Minimum \\
       $<0$&$>0$&$<0$ & negativ definit & Maximum\\
       $\geq0$&$\geq0$&$\geq0$& positiv semidefinit & Minimum o. Sattelpunkt \\
       $\leq0$&$\leq0$&$\leq0$& negativ semidefinit & Maximum o. Sattelpunkt \\
       sonst & & & indefinit & Sattelpunkt
    \end{tabular}
    \caption{Caption}
    \label{tab:my_label}
\end{table}


\section{Integrale bestimmen}
\subsection{Sollte man wissen}
\begin{align}
    \int \ln(x) dx &= \ln(x) \cdot x -x
\end{align}
\subsection{Substitution}
\begin{enumerate}
    \item Wähle Teilformel $g(x)$, setze $u = g(x)$
    \begin{itemize}
        \item Berechne $g(x)'$
        \item Schreibe $\frac{du}{dx} = g(x)'$, also $\frac{du}{g(x)'} = dx$
        \item Ersetze im Integral $g(x)$ durch $u$ und $dx$ durch $\frac{1}{g(x)'} du$
    \end{itemize}
    \item Falls noch $x$ auftaucht: Löse $u= g(x)$ nach $h(u) = x$ auf und ersetze jedes $x$ durch $h(u)$
    \item Berechne Integral bezüglich $u$
    \item Rücksubstitution
\end{enumerate}







\end{document}

